\documentclass{article}
\usepackage{amsmath, amsthm, amssymb}

\title{Proof of the Yang-Mills Gap Theorem}
\author{}
\date{}

\begin{document}

\maketitle

\section*{Chapter 2: Instantons and Tunneling in Yang-Mills Theory}

\subsection*{2.1 Introduction to Instantons}
Instantons are solutions to the Yang-Mills field equations in Euclidean space that represent tunneling events between different vacuum states. First introduced by Belavin, Polyakov, Schwarz, and Tyupkin in 1975, instantons have become a central tool in understanding non-perturbative effects in quantum field theory. Their importance lies in their ability to connect distinct vacua, thus providing a mechanism for phenomena such as quantum tunneling and the generation of a mass gap in Yang-Mills theory.

\subsection*{2.2 Mathematical Formulation}
The instanton solution in Euclidean Yang-Mills theory is obtained by considering the field strength tensor \( F_{\mu\nu}^a \) in the context of a gauge group \( SU(2) \). The field equations in Euclidean space are:
\[
D_\mu F^{\mu\nu} = 0
\]
where \( D_\mu \) is the covariant derivative. Instantons are self-dual solutions, satisfying:
\[
F_{\mu\nu} = \tilde{F}_{\mu\nu} = \frac{1}{2} \epsilon_{\mu\nu\rho\sigma} F^{\rho\sigma}
\]
This condition ensures that the Euclidean action \( S_E \) is minimized:
\[
S_E = \frac{1}{2g^2} \int d^4x \, F_{\mu\nu}^a F^{\mu\nu a} = \frac{8\pi^2 n}{g^2}
\]
where \( n \) is the topological charge, representing the winding number of the gauge field configuration.

\subsection*{2.3 Topological Charge and Energy Quantization}
The topological charge \( n \) is an integer that classifies the instanton solutions according to how many times they wrap around the vacuum manifold. This quantization leads to discrete energy levels associated with different instanton configurations:
\[
E_{\text{instanton}} = \frac{8\pi^2 n}{g^2}
\]
Since \( n \) is non-zero for non-trivial instanton configurations, the energy of the system is always positive, thereby preventing the existence of massless states. This quantization is key to establishing the mass gap in Yang-Mills theory.

\subsection*{2.4 Implications for the Mass Gap}
The existence of a positive energy contribution from instantons directly implies a positive mass gap in the theory. As instantons enforce energy quantization, they ensure that the vacuum cannot support excitations with zero energy. Therefore, the smallest possible energy of a non-trivial state is strictly positive, confirming the existence of a mass gap.

\end{document}